\section*{Introduction}

High-dimensional biological datasets are challenging for linear classifiers due to the small-sample, many-feature
regime and strong feature redundancy. Layered sparse classifiers provide an interpretable approach by learning
sequential neurons and removing selected features after each step.
This setting is also a classical use case for sparse linear modeling and embedded feature selection
\cite{tibshirani1996lasso,ng2004l1l2,saeys2007review}.

This work focuses on two complementary benchmarks. The first is the \texttt{colon\_cancer} microarray dataset,
one of the canonical benchmarks in bioinformatics and machine learning for high-dimensional classification
\cite{alon1999,guyon2002,dudoit2002,li2004}. It originates from the colon tumor vs.\ normal tissue profiling
study published in 1999 \cite{alon1999}. Across more than two decades, this benchmark has been repeatedly used
to evaluate feature selection, sparse linear models, and stability of classification under strong
$p \gg n$ conditions \cite{guyon2002,dudoit2002,li2004}. In comparative studies, \texttt{colon\_cancer}
typically yields high-80\% to mid-90\% accuracy ranges depending on validation protocol and feature-selection
strategy \cite{guyon2002,dudoit2002,li2004}.

The second benchmark is synthetic and is introduced for a different purpose. Real biological datasets provide
realistic distributional complexity, but they do not reveal the true predictive contribution of each feature.
As a consequence, they are suitable for measuring predictive performance and sparsity, but not for directly
measuring whether a method recovers the genuinely informative variables before drifting toward redundant or
purely noisy ones. This distinction is particularly important in high-dimensional gene-expression studies,
where feature selection and classifier evaluation can become confounded by selection bias if the protocol is
not carefully separated across training and test data \cite{ambroise2002,krawczuk2016feature}. To address
this limitation, we add a synthetic binary benchmark with known informative
features, controlled redundant features, and a predefined ranking of informative-feature strengths.
This perspective is also aligned with recent work emphasizing that feature-selection stability should be
treated as an important evaluation dimension in high-dimensional classification, not merely as a secondary
diagnostic after predictive accuracy \cite{lukaszuk2024importance}.

We compare three methods under the same protocol:
\texttt{cpl}, \texttt{svm}, and \texttt{logreg}. The goal is to assess how predictive performance, sparsity,
computational cost, and feature-recovery behaviour evolve across successive neurons in a controlled repeated
cross-validation setup.
The compared learners are all grounded in well-established sparse linear classification frameworks, including
L1-regularized logistic models and large-scale linear classifiers \cite{friedman2010glmnet,fan2008liblinear}.

Rather than assuming that aggregating many neurons must monotonically improve classification, we study a more
basic and more informative question: what does each next neuron contribute? In particular, we ask how much
predictive signal is captured by the first neurons, when later neurons still add complementary information, and
when they become dominated by redundant or noisy features.

Main contributions:
\begin{enumerate}[noitemsep]
\item A unified complex-layer framework evaluated with three linear sparse learners.
\item A repeated cross-validation benchmark on the classical \texttt{colon\_cancer} microarray dataset.
\item A synthetic benchmark with known feature relevance and known ranking of signal strengths.
\item Reproducible reporting of neuron-by-neuron accuracy, sparsity, weighted voting, and feature-recovery
diagnostics.
\end{enumerate}
